\section{Polyhedral degeneration and quantum volume asymptotics}
\label{sec:quantumdegeneration}

\subsection{Three different geometric quantities}

The connection to the volume conjecture becomes precise only after
separating three quantities:
\[
 A=\operatorname{Area}(\partial P),\qquad
 W(P)=\operatorname{Vol}_{\HH^3}(P),\qquad
 V_{\WP}=\int_{\M}\frac{\omega_{\WP}^d}{d!}.
\]
The first is a surface area, the second an enclosed three-dimensional
volume, and the third a measure of a space of shapes. The ordinary
knot volume conjecture concerns three-dimensional volume. For a
hyperbolic knot, in the reduced colored-Jones normalization
$J_N(\text{unknot};q)=1$, its familiar form is
\begin{equation}
 \lim_{N\to\infty}\frac{2\pi}{N}
 \log\left|J_N(K;e^{2\pi i/N})\right|
 =\operatorname{Vol}(S^3\setminus K).
 \label{eq:volumeconjecture}
\end{equation}
This remains a conjecture in general; see the colored-Jones
formulation in~\cite{MurakamiMurakami} and the tetrahedral
results in~\cite{Costantino}.

Suppose a shrinking family $P_A$, after lengths are multiplied by
$A^{-1/2}$, converges to a \emph{nondegenerate} Euclidean polyhedron
$P_{\mathrm E}$ of surface area one. Ordinary volume scaling gives
\begin{equation}
 W(P_A)=A^{3/2}W_{\mathrm E}(P_{\mathrm E})+o(A^{3/2}).
 \label{eq:individualflatvolume}
\end{equation}
In contrast, the corresponding moduli volume satisfies
\[
 V_{\WP}(A)=2^{-d}A^d V_{\mathrm{CH}}(\mu)+O(A^{d+1}).
\]
The different powers reflect different dimensions. Flattened
Euclidean limits and changes of polyhedral combinatorics require
additional care; they are not covered by the nondegeneracy
hypothesis in~\eqref{eq:individualflatvolume}.

\subsection{An explicit small-polyhedron expansion}

A controlled model gives more than the leading power. Let $Q$ be a
bounded nondegenerate convex Euclidean polyhedron, and place $P_r=rQ$ in
the curvature-$-1$ Klein ball. Its Euclidean planar faces are
hyperbolic geodesic faces. The Klein volume density is
\[
 \dd\operatorname{Vol}_{\HH^3}=(1-|x|^2)^{-2}\,\dd^3x.
\]
Consequently, whenever $r\sup_Q|y|<1$,
\begin{align}
 W(rQ)
 &=r^3\int_Q(1-r^2|y|^2)^{-2}\,\dd^3y\notag\\
 &=\sum_{k=0}^{\infty}(k+1)r^{3+2k}
                         \int_Q|y|^{2k}\,\dd^3y\notag\\
 &=r^3 W_{\mathrm E}(Q)
      +2r^5\int_Q|y|^2\,\dd^3y
      +3r^7\int_Q|y|^4\,\dd^3y+\cdots.
 \label{eq:Kleinvolumeexpansion}
\end{align}
This convergent expansion is elementary and exact for the specified
homothetic Klein path. It is not generally the fixed-conformal
angle path~\eqref{eq:anglepath}. Nor does it imply a Taylor expansion
of this form at ideal, flattening, or otherwise uncontrolled
degenerations. Unlike~\eqref{eq:exactWPvolumeflat}, it is an infinite
series: the finite polynomial for the moduli volume comes from a
quadratic cohomology class and a top intersection.

\subsection{Schl\"afli, doubling, and degeneration/regeneration}

For a compact hyperbolic polyhedron with interior dihedral angles
$\alpha_e$ and edge lengths $\ell_e$, Schl\"afli's formula is
\begin{equation}
 \dd W=-\frac12\sum_e\ell_e\,\dd\alpha_e.
 \label{eq:Schlafli}
\end{equation}
For the Klein family above, write
\[
 \ell_e(r)=rL_e+O(r^3),\qquad
 \alpha_e(r)=\alpha_e^{\mathrm E}+r^2b_e+O(r^4).
\]
The leading term of~\eqref{eq:Kleinvolumeexpansion} and
\eqref{eq:Schlafli} imply
\begin{equation}
 \boxed{\quad\sum_eL_e b_e=-3W_{\mathrm E}(Q).\quad}
 \label{eq:firstcurvaturevolume}
\end{equation}
Indeed, $\dd W/\dd r=3r^2W_{\mathrm E}+O(r^4)$, while
the right side of~\eqref{eq:Schlafli}, divided by $\dd r$, is
$-r^2\sum_eL_e b_e+O(r^4)$. Thus the first curvature correction
to the Euclidean dihedral angles already determines the Euclidean
enclosed volume. The sign reverses if exterior dihedral angles are
used. The same identity holds with
$\ell_e=A^{1/2}L_e+O(A^{3/2})$ and
$\alpha_e=\alpha_e^{\mathrm E}+A b_e+O(A^2)$ for a smooth
unit-area Euclidean limiting family.

Doubling a convex polyhedron across its faces produces a cone
three-manifold with singular graph given by its edges. The cone
angle along an edge is $2\alpha_e$, and its volume is $2W(P)$.
A shrinking polyhedron therefore supplies a direct example of
degeneration to a Euclidean cone three-manifold after rescaling.
This is the geometric setting suggested by Hodgson's thesis,
\emph{Degeneration and regeneration of geometric structures on
3-manifolds}~\cite{HodgsonThesis}. Later precise regeneration results
include Porti--Weiss~\cite{PortiWeiss}; their assumptions on cone
angles and the limiting structure must be checked before using
them for a particular doubled polyhedron. No unrestricted
regeneration theorem is being asserted here. The edge cone angles
$2\alpha_e$ are also distinct from the intrinsic vertex openings
$\theta_i$ of the boundary surface.

Rivin's \emph{Volumes of degenerating polyhedra---on a conjecture
of J.~W. Milnor} is directly relevant~\cite{RivinDegeneration}.
It studies continuous extension of volume toward degenerate angle
data, the Euclidean vanishing boundary, and quantitative control
through Schl\"afli's formula and polar geometry. Its estimates
allow logarithmically growing volume derivatives at degenerations;
continuity does not imply a smooth power series at every boundary
point. The controlled small-polyhedron expansion above gives a
particularly explicit local calculation on the Euclidean side of
this broader degeneration picture.

\subsection{The established tetrahedral quantum connection}

Quantum $6j$-symbols provide a concrete bridge to three-dimensional
geometry. Costantino obtains hyperbolic-volume exponential growth
for appropriate root-of-unity label regimes associated to truncated
tetrahedra~\cite{Costantino}. Chen--Murakami give, for a fixed
hyperbolic tetrahedron with at least one ideal or ultraideal vertex
and their prescribed admissible label sequences, an asymptotic of
the form
\begin{equation}
 |\{6j\}_r|\sim C\,r^{-3/2}(-\det G)^{-1/4}
                 \exp\left(\frac{r}{2\pi}W(T)\right),
 \label{eq:quantum6jstructure}
\end{equation}
where $G$ is the dihedral-angle Gram matrix and the prefactor and
labels depend on the stated conventions~\cite{ChenMurakami}.
In that regime the exponential growth rate is the tetrahedron's
hyperbolic volume. Different $6j$ normalizations and doubled
tetrahedral blocks change factors; the formula is not meant to
identify all such conventions.

At the Euclidean end, Roberts proves the Ponzano--Regge asymptotic
for classical $6j$-symbols: Euclidean tetrahedral edge lengths and
dihedral angles enter an oscillatory Regge phase, and Euclidean
volume enters its amplitude~\cite{Roberts}. This is an established
appearance of Euclidean tetrahedral geometry in quantum
asymptotics. It does not automatically identify the classical
large-label limit with a chosen root-of-unity hyperbolic limit.

In particular, the stated Chen--Murakami theorem does not cover
an arbitrary family of compact tetrahedra shrinking to Euclidean
ones. At a Euclidean tetrahedron the angle Gram determinant
vanishes, so the displayed fixed-shape prefactor itself warns of
nonuniformity. A continuous classical volume limit, even with
good quantitative estimates, does not by itself control the
quantum Hessian, the prefactor, or the error term in a joint limit.

\subsection{A proposed joint scale and a model integral}

Suppose a suitable quantum object, in a specified geometric branch
and choice of variables, has an exponential contribution
$\exp(W(P_A)/\hbar)$. Equation~\eqref{eq:individualflatvolume}
then singles out the dimensionless ratio
\begin{equation}
 \lambda=\frac{A^{3/2}}{\hbar},\qquad
 A\asymp\hbar^{2/3}.
 \label{eq:jointscale}
\end{equation}
This is the regime in which geometric degeneration competes
with the semiclassical limit. At this stage it is a scaling prediction;
Section~\ref{sec:quantumAiry} proves a uniform Airy statement for an
explicitly normalized regular $b$-$6j$ family. A different
polarization may involve a Regge or Legendre-transformed action,
so the quantum object and phase must be specified first.

There is a useful classical model making the potential role of
 the WP--Thurston limit explicit. Let $K\subset\M_{0,n}$ be a
 fixed compact subset with controlled nondegenerate Euclidean
 realizations, and define
\begin{equation}
 \mathcal Z_K(A,\hbar)=\int_K
       \exp\left(\frac{W(P_A(X))}{\hbar}\right)
       \frac{\omega_{\WP,A}^d}{d!}.
 \label{eq:modelpartition}
\end{equation}
This is a constructed geometric integral, not an identification
with a colored-Jones or other known quantum invariant. Under the
corresponding uniform geometric convergence on $K$, take
$A\to0$, $\hbar\to0$, and $A^{3/2}/\hbar\to\lambda_0\in(0,\infty)$.
The local limits then give
\begin{equation}
 2^d A^{-d}\mathcal Z_K(A,\hbar)\longrightarrow
 \int_K e^{\lambda_0 W_{\mathrm E}(P_{\mathrm E}(X))}
                 \frac{\omega_{\mathrm{CH},\mu}^d}{d!}.
 \label{eq:modelpartitionlimit}
\end{equation}
Extending this formula to the full noncompact space requires
 domination and control of collision and flattened strata.
 Identifying its measure and phase with those of an actual quantum
 state integral is a further mathematical problem.

This separates the two potentially compatible roles: enclosed
three-dimensional volume supplies an action, while the WP form
supplies a measure on boundary shapes. Their Euclidean limits
are respectively Euclidean polyhedral volume and the
Thurston--Deligne--Mostow measure. Establishing that both occur
together in a particular quantum construction would provide the
substantive connection to the volume conjecture.


